Section \ref{sec.stability} describes the process in the mean.
The mean is not enough to design a measurement.
If we intend to read the state from a window sum of recent events,
what decides how long the window must be is not the size of the intensity
but the size of its fluctuations:
a window is long enough when its sum is dominated by the signal rather than by the noise of
counting.
This section computes the stationary second moments of the state,
in closed form, for the exponential kernel.

\subsection{The generator}

By \eqref{eq.stateDynamics} the pair $(\multiCountingProc, \decayedCounts)$ is a
piecewise-deterministic Markov process,
and for $f$ continuously differentiable the process
\begin{equation*}
 f(\decayedCounts_t) - f(\decayedCounts_0)
 - \int_{0}^{t} \left[
 -\decay \, \decayedCounts_u \cdot \gradient f(\decayedCounts_u)
 + \sum_{e} \intensity(u) \Big( f(\decayedCounts_u + e_e) - f(\decayedCounts_u) \Big)
 \right] du
\end{equation*}
is a local martingale.
Applying this with $f(s) = s$ and with $f(s) = s s^{\transpose}$,
and exchanging expectation and derivative --- which is legitimate as soon as
$\Expectation \vert \decayedCounts \vert^2 < \infty$,
and we shall see that this holds exactly when $\branchingRatio<1$ ---
gives the two equations of this section.

\subsection{The mean state}

With $f(s)=s$,
\begin{equation*}
 \frac{d}{dt}\Expectation[\decayedCounts_t]
 = -\decay \Expectation[\decayedCounts_t] + \Expectation[\intensity[ ](t)]
 = \baseIntensity + (\excitation - \decay I) \Expectation[\decayedCounts_t] ,
\end{equation*}
using \eqref{eq.intensityFromState}.
At stationarity the left-hand side vanishes, so writing
$m := \Expectation[\decayedCounts]$ we have
$(\excitation - \decay I) m = -\baseIntensity$, and
\begin{equation}\label{eq.meanState}
 m = \frac{\stationaryIntensity}{\decay} .
\end{equation}
Indeed $\decay m = \baseIntensity + \excitation m$ is \eqref{eq.intensityFromState} in the mean,
which is \eqref{eq.stationaryIntensity}.
The identity
\begin{equation}\label{eq.meanStateIdentity}
 (\excitation - \decay I) \, \frac{\stationaryIntensity}{\decay} = -\baseIntensity
\end{equation}
is the one that does the work below.

\subsection{The covariance}

With $f(s) = ss^{\transpose}$ the jump term needs care.
At an event of type $e$ the state gains the unit vector $e_e$,
so $\decayedCounts \decayedCounts^{\transpose}$ gains
$\decayedCounts e_e^{\transpose} + e_e \decayedCounts^{\transpose} + e_e e_e^{\transpose}$.
Two facts about the jumps enter here and both are hypotheses, not conveniences:
the jumps of $\decayedCounts$ are \emph{unit} vectors,
which is what makes the last term $e_e e_e^{\transpose}$;
and no two components jump together, by \eqref{eq.noCommonJumps},
which is what keeps that term diagonal.
Writing $M := \Expectation[\decayedCounts\decayedCounts^{\transpose}]$ and using
$\Expectation[\decayedCounts \intensity[ ]^{\transpose}]
= m\baseIntensity^{\transpose} + M\excitation^{\transpose}$,
stationarity gives
\begin{equation*}
 0 = (\excitation - \decay I) M + M (\excitation - \decay I)^{\transpose}
 + \baseIntensity m^{\transpose} + m \baseIntensity^{\transpose}
 + \diag(\stationaryIntensity).
\end{equation*}
Substituting $M = V + m m^{\transpose}$, where $V := \CoVariance(\decayedCounts)$,
the cross terms cancel exactly, since by \eqref{eq.meanStateIdentity}
\begin{equation*}
 (\excitation - \decay I) m m^{\transpose} + m m^{\transpose}(\excitation - \decay I)^{\transpose}
 = -\baseIntensity m^{\transpose} - m\baseIntensity^{\transpose} .
\end{equation*}
What remains is a Lyapunov equation.

\begin{prop}\label{prop.lyapunov}
 Let $\branchingRatio < 1$.
 The stationary covariance $V$ of the state $\decayedCounts$ is the unique solution of
 \begin{equation}\label{eq.lyapunov}
  (\excitation - \decay I) V + V (\excitation - \decay I)^{\transpose}
  + \diag(\stationaryIntensity) = 0 ,
 \end{equation}
 and the stationary covariance of the intensity is
 $\CoVariance(\intensity[ ]) = \excitation V \excitation^{\transpose}$.
\end{prop}

The forcing term is the diagonal matrix of jump intensities:
fluctuation is injected at the rate at which events occur,
one unit of variance per event, and no covariance between types because no two fire together.
The homogeneous part transports it,
and the whole content of the equation is the balance between the two.

Uniqueness is worth stating separately, because it is the cleanest statement of what
subcriticality buys at second order.
A Lyapunov equation $NV + VN^{\transpose} + Q = 0$ has a unique solution exactly when $N$ is
Hurwitz, and $\excitation - \decay I$ is Hurwitz if and only if
$\branchingRatio(\excitation) < \decay$, that is, if and only if $\branchingRatio < 1$.
Subcriticality is not merely sufficient for bounded second moments: it is equivalent to them.

\subsection{The scalar case, and when it may be used}

Contracting \eqref{eq.lyapunov} with $\mathbf{1}^{\transpose}$ on the left and
$\mathbf{1}$ on the right gives a scalar equation as soon as
$\mathbf{1}^{\transpose}\excitation$ is proportional to $\mathbf{1}^{\transpose}$,
that is, as soon as $\excitation$ has constant column sums.
Those sums are then necessarily $\decay\branchingRatio$, and
$\mathbf{1}^{\transpose}(\excitation - \decay I) = -\decay(1-\branchingRatio)\mathbf{1}^{\transpose}$,
so that
\begin{equation}\label{eq.scalarVariance}
 \Variance \big( \mathbf{1}^{\transpose}\decayedCounts \big)
 = \frac{\totalRate}{2\decay(1-\branchingRatio)} .
\end{equation}
In one dimension this is $\Variance(\decayedCounts) = \stationaryIntensity/(2\decay(1-\branchingRatio))$,
and it is the formula usually quoted.

Three cautions attach to it.
It is a statement about $\mathbf{1}^{\transpose}\decayedCounts$,
the total decayed count,
and not about any single component;
it is not a statement about the total intensity, which in the same special case has
variance $\decay\branchingRatio^{2}\totalRate/(2(1-\branchingRatio))$;
and it requires the constant column sums, which are not generic.
On the parametrisation of these notes the column sums are not constant,
and \eqref{eq.scalarVariance} understates $\mathbf{1}^{\transpose}V\mathbf{1}$ by $17.7\%$
--- a discrepancy that does not depend on $\decay$, since both sides scale as $1/\decay$.
How large it is depends entirely on how far from constant the column sums are, which is a
property of the kernel and not something the formula can warn about.
When the covariance is wanted, \eqref{eq.lyapunov} is a six-by-six linear solve.
